\section{Evaluation}
\label{sec:evaluation}

\subsection{RQ1: cJSON, baseline versus LLM-refined}

Table~\ref{tab:cjson-comparison} reports the five-seeded-run-per-arm
comparison described in Section~\ref{sec:methodology}. Because every one of
the five refined-arm runs exceeds every one of the five baseline-arm runs on
both acceptance rate and structural fingerprints, and every refined-arm run
falls below every baseline-arm run on rejection signatures, all three
metrics exhibit complete rank separation between arms. For a two-sample
comparison with five observations per arm and no ties, complete separation
is the most extreme outcome an exact (permutation-based) Mann-Whitney U test
can observe: $U{=}25$ (or, equivalently, $U{=}0$ from the other arm's
perspective) out of a maximum of 25, giving an exact two-sided
$p = 2/\binom{10}{5} \approx 0.0079$ for each of the three metrics. We report
this as the exact result it is, not as evidence of an arbitrarily small
true $p$-value: with $n{=}5$ per arm, $p \approx 0.0079$ is the smallest
attainable exact $p$-value regardless of effect size, so it should be read
as ``the strongest result this sample size can show,'' not as unusually
strong evidence beyond what five runs per arm can support.

\begin{table}[t]
\centering
\caption{cJSON: baseline vs.\ LLM-refined, five seeded runs per arm (mean
over runs). Exact two-sided Mann-Whitney $p \approx 0.0079$ for all three
metrics (complete rank separation, $n{=}5$ per arm).}
\label{tab:cjson-comparison}
\begin{tabular}{lrrr}
\toprule
Metric & Baseline & Refined & Change \\
\midrule
Acceptance rate & 57.64\% & 96.77\% & $+$39.13 pp \\
Structural fingerprints (sum/run) & 276.0 & 616.2 & $+$123.3\% \\
Rejection signatures (sum/run) & 92.4 & 29.6 & $-$68.0\% \\
Crashes / timeouts / signals & 0 & 0 & --- \\
\bottomrule
\end{tabular}
\end{table}

\textbf{Unequal budgets.} Baseline runs execute 500 inputs on average;
refined runs execute 1700, because the refinement loop ran more iterations
per run on average. Every count-based quantity (accepted, rejected,
fingerprints, signatures as raw totals rather than rates) therefore scales
with a budget that differs between arms and is not directly comparable;
acceptance rate is scale-free and unaffected by this. We report the
budget-sensitive metrics above as \emph{sums over a run's iterations}
specifically so the comparison is transparent about this rather than implying
parity.

\textbf{No crashes in either arm.} Zero sanitizer crashes, timeouts, or
fatal signals occurred in either arm across all ten runs. The effect
measured here is therefore on generator quality -- validity and structural
diversity -- not on vulnerability discovery in cJSON; we do not claim or
imply that refinement found, or was likely to find, a memory-safety bug in
this target.

\subsection{RQ2: generalization to parson}

We reused every pipeline component unchanged against a second, independent
target, parson (commit \texttt{ba29f4e}, v1.5.3), with a human standing in
for the LLM proposer call and every resulting proposal passing through the
identical validation and execution path (Section~\ref{sec:methodology}).
Reconciling the same JSON grammar against parson's real behavior surfaced
three adaptations, two of which go the \emph{opposite} direction from the
cJSON adaptations reported in Section~\ref{sec:methodology} under the
identical grammar: (i) unlike the cJSON harness, parson's parse entry point
does not require full-buffer consumption, so trailing garbage after a valid
value is silently accepted -- a superset of what the grammar's
\texttt{json : value EOF ;} production specifies; (ii) parson explicitly
rejects any object containing a duplicate key, propagating the rejection up
as a failure for the whole object -- a subset of what cJSON accepts, even
though the grammar itself is silent on duplicate keys for both; (iii) parson
enforces a fixed nesting limit (\texttt{MAX\_NESTING = 2048}), rejecting
deeper structures outright rather than exhibiting unbounded recursion. That
the same formal grammar under-determines a directly opposite validity
decision (duplicate keys: accept vs.\ reject) on two independently
implemented parsers is, on its own, evidence that a purely formal grammar
specification cannot supply the target-specific signal needed to fuzz either
implementation precisely.

\begin{table}[t]
\centering
\caption{parson: five real refinement iterations, 500 examples each.}
\label{tab:parson-iterations}
\small
\begin{tabular}{lrrr}
\toprule
Iter.\ & Focus (abridged) & Accepted/Rejected & Fingerprints \\
\midrule
1 & Grammar-seeded generator: recursion, numeric edge cases, duplicate keys, trailing garbage & 285 / 215 & 206/500 \\
2 & Wide objects/arrays (up to 4000 keys) to hit rehash/resize paths unreached by iter.\ 1 & 341 / 159 & 273/500 \\
3 & Hash-bucket collisions, combined wide+deep structures, surrogate escapes at closing quote & 442 / 58 & 254/500 \\
4 & Internal error-cleanup free path: structures valid until the final token & 85 / 415 & 214/500 \\
5 & Combines iter.\ 4's free-path shapes with a sibling-cascade free case & 139 / 361 & 302/500 \\
\bottomrule
\end{tabular}
\end{table}

\textbf{Result.} No crash was found across all 3{,}000 sanitizer-instrumented
executions (500 baseline plus $5 \times 500$ refinement iterations), verified
by grepping every run's stderr for AddressSanitizer, UndefinedBehaviorSanitizer,
and LeakSanitizer text and checking for any non-\texttt{accepted}/\texttt{rejected}
status. Manual review of the highest-risk code paths -- \texttt{parse\_utf16}'s
raw pointer walk across surrogate pairs, \texttt{process\_string}'s
output-buffer sizing, \texttt{json\_object\_grow\_and\_rehash}'s linear
probing, and, informed specifically by iterations 4--5, \texttt{json\_object\_add}'s
duplicate-key failure branch and \texttt{parse\_object\_value}/\texttt{parse\_array\_value}'s
missing-bracket failure branch -- found that each relies on the harness's own
null terminator as a safe, in-bounds stopping condition, and that internal
cleanup calls to \texttt{json\_value\_free} are consistent with the harness's
own top-level call (an object's key count is never incremented before a
duplicate is detected, so cleanup never iterates past the last successfully
added slot). We report this as a defensible, evidence-supported explanation
for the absence of a crash within this budget, not as proof that none
exists: five iterations against one pinned commit, exercising the parse and
internal free paths but not parson's public mutation API
(\texttt{json\_object\_remove}, \texttt{json\_array\_remove}, and the
\texttt{\_replace\_*} family), is a bounded search, and we do not claim more
than that bound supports.
